%% kapa: exam=1 round=30 year=2019 period=1 form=A subject=civil_law subject_ko=민법 seq=40 number=40
%% source: Q-Net 감정평가사 기출문제 (cst00309) artlSeq=5205205 file=extracted/1008/5205205/2019년 제30회 감정평가사 1차시험 문제지/2019년 제30회 감정평가사 1차시험 문제지/2019년도 제30회 감정평가사 1차 1교시 A형.hwp
%% conversion: kordoc 4.24.51 -> tools/gen1.py (auto); review=proofread-sonnet; flags=none
\begin{question}{40}
근저당권에 관한 설명으로 옳지 않은 것은? (다툼이 있으면 판례에 따름)
\begin{choices}
  \item 근저당권의 피담보채무는 원칙적으로 당사자가 약정한 존속기간이나 결산기가 도래한 때에 확정된다.
  \item 장래에 발생할 특정의 조건부채권을 피담보채권으로 하는 근저당권의 설정은 허용되지 않는다.
  \item 근저당부동산의 제3취득자는 피담보채무가 확정된 이후에 채권최고액의 범위 내에서 그 확정된 피담보채무를 변제하고 근저당권의 소멸을 청구할 수 있다.
  \item 근저당권자가 피담보채무의 불이행을 이유로 경매신청을 하여 경매 신청시에 근저당 채무액이 확정된 경우, 경매개시 결정 후 경매신청이 취하되더라도 채무확정의 효과가 번복되지 않는다.
  \item 채권최고액은 반드시 등기되어야 하지만, 근저당권의 존속기간은 필요적 등기사항이 아니다.
\end{choices}
\end{question}
\answer{2}
